\documentclass[11pt]{article}
\usepackage{amsmath,amssymb,amsthm,hyperref}
\newtheorem{theorem}{Theorem}
\newtheorem{lemma}[theorem]{Lemma}
\newtheorem{corollary}[theorem]{Corollary}
\newtheorem{proposition}[theorem]{Proposition}
\title{A pole obstruction to parity decoders on an open box}
\author{}
\date{November 22, 2026}
\begin{document}
\maketitle

This note concerns a proposed decoder in its real parameter variables.
It is not a lower bound for Boolean head complexity: its sign requirement
holds throughout an open box, not just at Boolean vertices or at one
chosen representative of each permutation. No proof of the requested
universal upper bound is claimed here.

\begin{lemma}[Edges along coordinate fibers]
Let $T$ be a finite nonempty subset of $\mathbb R^r$, with $N=|T|$.
In each line parallel to a coordinate axis, join consecutive points of
$T$ on that line. The resulting simple graph has at most
$N\log_2 N/2$ edges. A vertex having another point of $T$ in each of
$d$ distinct coordinate directions has degree at least $d$.
\end{lemma}
\begin{proof}
The degree assertion follows by taking a neighbor in each of the stated
fibers. These neighbors are distinct. For the edge bound, induct on $N$.
For $N=1$ it is immediate. Otherwise some coordinate is nonconstant on
$T$. Cut between two consecutive values of this coordinate, partitioning
$T$ into nonempty sets of sizes $a,b$. Edges crossing the cut form a
matching: only edges in the cut direction can cross, and each fiber
has at most one such edge. There are therefore at most $\min(a,b)$
crossing edges. Within each part the edges are precisely those of its
own consecutive-point graph. Induction gives
\[
 |E(T)|\le\frac{a\log_2a+b\log_2b}{2}+\min(a,b).
\]
For $p=a/(a+b)$, the binary entropy
$h(p)=-p\log_2p-(1-p)\log_2(1-p)$ satisfies
$h(p)\ge2\min(p,1-p)$. This follows from concavity on $[0,1/2]$
with endpoint values $0,1$, followed by symmetry. Substitution gives
the required upper bound.
\end{proof}

\begin{theorem}[An exponential obstruction for continuous parity]
Let $r\ge2$ and let $U$ be an open box containing the origin in
$\mathbb R^r$. Suppose
\[
 S(z)=P_0(z)+\sum_{h=1}^H\frac{A_h(z)}{B_h(z)},
\]
where $P_0$ and all $A_h,B_h$ are real affine forms and every $B_h$ is positive
on $U$. Suppose that, whenever every coordinate of $z\in U$ is nonzero,
\[
 \operatorname{sign}S(z)=\operatorname{sign}\prod_{i=1}^r z_i.
\]
Then $H\ge2^{r-2}$. More precisely, after the reductions in the proof,
the number $K$ of distinct nonconstant genuine poles satisfies
\[
 K\ge r-1,\qquad K(r-1)\le(K+1)\log_2(K+1).
\]
\end{theorem}
\begin{proof}
Normalize every denominator to have constant term one, which is possible
because $B_h(0)>0$. Combine equal normalized denominators. Absorb
constant-denominator terms into an affine polynomial $P$, and absorb any
numerator proportional to its denominator into the constant part of $P$.
After deleting zero terms, write
\[
 S=P+\sum_{h=1}^K\frac{A_h}{1+b_h\cdot z}.
\]
Here the $b_h$ are distinct and nonzero, and $A_h$ is not a scalar
multiple of $1+b_h\cdot z$. We have $K\le H$.

Continuity and the two opposite signs on the sides of each coordinate
hyperplane imply $S=0$ on a relatively open part of that hyperplane.
Clearing denominators turns this into a polynomial identity on the
whole hyperplane. The cleared numerator is a nonzero polynomial of
degree at most $K+1$. Each coordinate variable divides it, and the
variables are distinct prime factors. Thus their product divides it,
so $K+1\ge r$. In particular $K\ge1$.

Put $T=\{0,b_1,\ldots,b_K\}$. Fix $h$ and a coordinate $i$ for which
there is no other point of $T$ differing from $b_h$ only in coordinate
$i$. The restriction of $1+b_h\cdot z$ to $z_i=0$ is nonconstant;
otherwise $b_h$ and $0$ would be such neighbors. This restriction is
also different from every other normalized restricted denominator;
equality would give the forbidden neighbor $b_j$.

In the polynomial ring of the remaining $r-1$ variables, this
nonconstant affine denominator is irreducible and prime. In the cleared
identity $S|_{z_i=0}=0$, every summand except the $h$th is divisible by
it. None of the other denominator factors is divisible by it, since
all have constant term one and none is equal to it. Consequently it
divides $A_h|_{z_i=0}$. Since this numerator is affine, there are scalars
$\lambda_i,\mu_i$ with
\[
 A_h(z)=\lambda_i(1+b_h\cdot z)+\mu_i z_i.
\]
The same conclusion for a different coordinate $j$ gives
$\lambda_j=\lambda_i=A_h(0)$ by evaluating at zero, and then
$\mu_i z_i=\mu_j z_j$ as polynomial identities. Since $i\ne j$, both
coefficients vanish. This would make $A_h$ proportional to its
denominator, contrary to the reduction. Thus each $b_h$ can lack a
neighbor in at most one coordinate direction.

Apply the preceding graph lemma to $T$. Each of its $K$ nonzero vertices
has degree at least $r-1$; no degree condition is imposed on the extra
vertex $0$. The degree sum and edge bound give
\[
 K(r-1)\le2|E(T)|\le(K+1)\log_2(K+1).
\]
Set $N=K+1$. We already proved $N\ge r$. Therefore
\[
 \log_2N\ge\frac{N-1}{N}(r-1)
 \ge\frac{(r-1)^2}{r}=r-2+\frac1r>r-2.
\]
Since $N$ is an integer and $2^{r-2}$ is an integer, this implies
$K=N-1\ge2^{r-2}$ and hence the asserted bound on $H$.
\end{proof}

\begin{proposition}[A construction within a factor two]
On any sufficiently small box about zero, continuous parity in $r$
variables has a representation by $2^{r-1}$ positive affine fractions.
\end{proposition}
\begin{proof}
Choose $t>0$ and a box on which $t\sum_{i<r}|z_i|<1$. Set
\[
 S(z)=z_r\sum_{\epsilon\in\{-1,1\}^{r-1}}
       \frac{\prod_{i<r}\epsilon_i}
       {1-t\sum_{i<r}\epsilon_i z_i}.
\]
All denominators are positive. Using the convergent Laplace integral for
each reciprocal and summing the finitely many terms gives
\[
 S(z)=z_r\int_0^\infty e^{-s}
                \prod_{i<r}2\sinh(stz_i)\,ds.
\]
When all coordinates are nonzero, the integrand for $s>0$ has the
constant strict sign $\prod_{i<r}z_i$. The integral converges, since the
exponential growth rate of the product is less than one. Its product
with $z_r$ has exactly the required strict sign.
\end{proof}

\begin{corollary}[Order parity and matchings]
Let $E$ be a set of pairs of $p$ indices with a matching of size $m\ge2$.
Suppose a decoder formed from $H$ affine fractions in
$\theta\in I^p$, with all denominators positive on this open box, has
the sign of
\[
 \prod_{\{i,j\}\in E,\ i<j}(\theta_j-\theta_i)
\]
whenever the coordinates are distinct. Then $H\ge2^{m-2}$.
\end{corollary}
\begin{proof}
Assign a distinct center in $I$ to each matched pair and to every
unmatched index. For the two indices in the $k$th matched pair use
$\theta_i=c_k-z_k/2$ and $\theta_j=c_k+z_k/2$, in their increasing index
order. On a sufficiently small box in $z$, all factors outside the
matching stay nonzero with constant signs, and the matched factors are
exactly $z_1,\ldots,z_m$. All coordinates of $\theta$ are distinct
whenever the $z_k$ are nonzero. Restriction of the decoder gives affine
numerators and positive affine denominators on this small box. Multiply
its score by the fixed sign of the other factors and apply the theorem.
\end{proof}

\paragraph{Scope of the obstruction.}
The order-parity decoder must work on every distinct-coordinate point
of the open box, and its coefficients must depend affinely on the
stored scalars so that restriction still gives affine forms. A decoder
that works only at specially chosen representatives of the permutation
orders, or assigns its head parameters nonlinearly from the stored
scalars, is outside this result. The bound therefore blocks the stated
continuous decoder construction, not the full conjecture about $W(n)$.

\end{document}
